The speed of a GPU kernel is usually decided not by the arithmetic it performs, but by how that arithmetic is laid out on the hardware. The same matrix multiplication can run several times faster or several times slower depending on how its work is tiled across thread blocks, how many warps are launched per block, and how aggressively the compiler overlaps memory and compute. Triton exposes these choices directly, which is convenient for an expert and expensive for everyone else: a developer rarely has the time to benchmark every plausible configuration, so they have to decide up front which kinds of evidence are worth paying for before running the kernel at full cost.

This paper studies that trade-off under a fixed measurement budget. We ask whether cheap signals that can be read immediately after compilation, such as register usage, shared-memory footprint, and occupancy estimates, together with a small amount of runtime profiling, can pick a better schedule than default settings or an equally budgeted random search. To answer this, we build \system, a framework that treats candidate generation, admissibility checks, compile-time signals, limited profiling, and held-out runtime evaluation as distinct tiers of an evidence hierarchy, and we evaluate it on matrix multiplication and layer normalization kernels running on a homogeneous pool of NVIDIA RTX A6000 GPUs. Throughout the paper, we quote results as a geometric-mean speedup: a value of $1.03\times$ means the tuned kernel finishes in about $97\%$ of the default runtime on the same inputs, while a value of $0.83\times$ means it is roughly $20\%$ \emph{slower} than the default.

Our findings are mixed, which we think is the useful part. Cheap compile-time signals turn out to be strong pruning tools and weak rankers: they reject obviously poor schedules decisively, yet they struggle to separate the plausible ones from each other. The choice of workload changes the conclusion: the same selector looks noticeably stronger on conveniently aligned shapes than on representative ones with irregular aspect ratios. Runtime profiling helps on some kernels and regimes and hurts on others, and one reduction-splitting knob that looked promising on a narrow search space collapsed once the space was enlarged, which is why we retire it from the final evaluation. After those retreats, a deliberately conservative selector still recovers a small but credible improvement on representative matrix multiplication: a geometric-mean speedup of $1.03\times$ over default settings, compared with $0.81\times$ for an earlier compile-signal-only ranking and $1.00\times$ for naive random search under the same matched budget. We keep the final number at the ``bounded improvement'' tier because it is stronger on repeated-run stability than on a cross-seed sweep, and because it survives the harsher tests above rather than replacing them. The contribution of the paper is therefore not the headline speedup itself, but the clearer account of what lightweight evidence can and cannot buy once the budget is fixed and the workloads are realistic.
